\begingroup
\renewcommand{\arraystretch}{1.45}
\fancyhead[C]{}
\clearpage\fancyhead[L]{\small BEMERKUNGEN.}\pdfbookmark[0]{Bemerkungen}{final-bemerkungen}\SourceLeaf{500}
\begin{center}\Large\textbf{BEMERKUNGEN.}\end{center}
\begin{center}\rule{.09\textwidth}{.4pt}\end{center}

In diesem dritten Band von \textsc{Gauss}’ Werken habe ich alle Abhandlungen und Aufsätze aus dem
Gebiete der allgemeinen Analysis und zwar speziell aus der Theorie der algebraischen Funktionen, der
\textsc{Gauss}’schen Reihen und der elliptischen Funktionen, sowie einige Mitteilungen das \textsc{Pfaff}sche Theorem
und die Konstruktion von Logarithmentafeln betreffend vereinigt. Sie bestehen aus einer Doktordissertation
(früher in Quart gedruckt); aus früher veröffentlichten Abhandlungen: fünf in den \textit{‘Commentationes societatis regiae scientiarum Gottingensis’} (Quart), einer in den \textit{‘Abhandlungen der Königlichen Gesellschaft der}
\textit{Wissenschaften zu Göttingen’} (Quart) und einer im \textsc{Crelle}schen \textit{‘Journal für reine und angewandte Mathematik’} (Quart); ferner aus sechs Anzeigen eigener Abhandlungen in den \textit{‘Göttingischen Gelehrten Anzeigen}
\textit{und Nachrichten’} (Oktav); aus zwölf Mitteilungen über nicht eigene Schriften in den \textit{‘Göttingischen gelehrten Anzeigen’} (Oktav) (von \textsc{Gauss} nicht unterzeichnet aber in Bezug auf seine Autorschaft durch die Akten
der Göttinger Universitäts-Bibliothek verifiziert), in der \textit{‘Monatlichen Correspondenz zur Beförderung der Erd-}
\textit{und Himmels-Kunde herausgegeben vom Freiherrn \textsc{von Zach}’} (Oktav), in den \textit{‘Astronomischen Nachrichten}
\textit{herausgegeben von \textsc{Schumacher}’} (Quart); und auch aus zwei Erläuterungen in der \textit{‘Monatlichen Correspondenz’}
und in \textsc{Vega}’s \textit{‘Sammlung von Hülfstafeln’} (Quart) zu den \textsc{Gauss}’schen Additions- und Subtraktions-Logarithmentafeln. Diese Tafeln selbst habe ich hier nicht abdrucken lassen, weil sie sehr verbreitet sind und
das Format dieser Werke zum Gebrauch solcher Tafeln unbequem sein würde.

Bei der Redaktion habe ich dieselben Grundsätze befolgt wie in den früheren Bänden. Zur besseren
Übersicht der Gegenstände in einem so umfangreichen Band sind die Hauptlehrsätze auf gleiche Weise
durch den Druck ausgezeichnet. Zum leichteren Gebrauch sowohl der älteren Ausgaben als der vorliegenden
habe ich bei den Verweisungen auf Schriften, die nicht in diesem Band selbst sich finden, statt des Ortes der
Veröffentlichungen die eigenen Titel sowie statt der Nummer der Seite die der Artikel angegeben. Auf
Seite 20 in Zeile 15 und 16 ist gemäß einer handschriftlichen Bemerkung \textit{‘objectionem secundam et quar\SourceLeaf{501}%
tam’} statt \textit{‘objectionem tertiam et quartam’} und in Zeile 19 ebenso \textit{‘objectionem tertiam’} statt \textit{‘objectionem}
\textit{primam’} gesetzt. Außerdem unterscheidet sich die vorliegende Ausgabe von den früheren derselben Schriften nur durch die Berichtigung einiger Druckfehler, wobei ich zum Teil die \textsc{Gauss}’schen Originalmanuskripte benutzen konnte.

Aus dem handschriftlichen Nachlass habe ich aufgenommen: die den Seiten 30 und 112 beigefügten Noten, den zweiten Teil der Abhandlung \textit{‘Disquisitiones generales circa seriem infinitam etc.’}, eine
ausführliche Abhandlung über Interpolation mit einigen in meinen Bemerkungen Seite 328 u.\,f. erläuterten
Zusätzen, und eine Reihe von Abhandlungen, die sich auf die elliptischen Funktionen beziehen.

Die handschriftlichen Aufzeichnungen habe ich hier wie auch früher bis auf Berichtigung von unerheblichen Schreibfehlern unverändert abdrucken lassen und meine Einschaltungen, die mit Ausnahme der
\textit{‘Fortsetzung der Untersuchungen über das arithmetisch-geometrische Mittel’} nur in kurzen Sätzen bestehen,
durch [Einklammerung] abgesondert. Die geschichtlichen Angaben und erforderlichen Zusätze für den Nachlass über hypergeometrische Reihen und Interpolation habe ich in den jenen Abhandlungen unmittelbar
folgenden Bemerkungen zusammengestellt.

\bigskip
Die Redaktion der \textsc{Gauss}’schen Arbeiten über elliptische Funktionen wurde seitens der Königlichen
Gesellschaft der Wissenschaften zu Göttingen von \textsc{Riemann} gewünscht und ihm zu diesem Zweck die Handschriften übergeben. Leider hat er weder eine schriftliche noch mündliche Mitteilung aus diesen seinen
Studien hinterlassen. Erst nach dem bedauernswerten allzufrühen Tode \textsc{Riemann}s und nachdem der vorliegende dritte Band bis auf jenen Teil gedruckt war, konnte ich die \textsc{Gauss}’schen Handschriften zu meiner Bearbeitung übernehmen.

\bigskip
\textsc{Gauss} hat von seinen Untersuchungen der Funktionen, die wir jetzt die elliptischen nennen, nur einen Teil veröffentlicht: eine Anwendung dieser Theorie auf die höhere Arithmetik in der \textit{Summatio quarundam serierum singularium} 1808 \textit{September} und eine Anwendung auf die Bestimmung der Säkularstörungen
der Planeten in der \textit{Determinatio attractionis, quam in punctum quodvis positionis datae exerceret planeta, si}
\textit{ejus massa per totam orbitam ratione temporis, quo singulae partes describuntur, uniformiter esset dipertita,} 1818
\textit{Januar.} Den größeren Teil hat er niedergelegt in einigen unvollständigen Entwürfen zu Anfängen verschiedener Abhandlungen und in zahlreichen zwischen andern Arbeiten sehr zerstreuten Aufzeichnungen einzelner
Formeln. Diese im handschriftlichen Nachlass befindlichen Untersuchungen habe ich hier in einzelnen
Gruppen zusammengestellt je nachdem sie vom Algorithmus des Arithmetisch-Geometrischen Mittels oder einem anderen diesem analogen und mit diesem in Verbindung gesetzten Algorithmus ausgehen oder aber
sich auf die speziellen lemniskatischen Funktionen beziehen oder endlich die Darstellungen der allgemeinen
Funktionen durch Produkte als wesentliches Hilfsmittel gebrauchen. Die Untersuchung des \textit{Pentagramma}
\textit{mirificum} bildet eine Anwendung der Fünfteilung der ganzen elliptischen Integrale.

\bigskip
Die Abhandlung mit den beiden Abschnitten \textit{de origine proprietatibusque generalibus numerorum mediorum arithmetico-geometricorum} und \textit{de functionibus transcendentibus quae ex differentiatione mediorum}
\providecommand{\TStart}{\hspace*{.5em}}\providecommand{\TCont}{\hspace*{1em}}\providecommand{\TD}{\hspace{.3em}.\hspace{.3em}.}
\SourceLeaf{502}
\textit{arithmetico-geometricorum oriuntur} bildet die erste und zwar eine sehr sorgfältig geschriebene Aufzeichnung in einem Handbuch, welches, wie der Titel besagt, von dem Jahre 1800 an benutzt ist. Die Beschäftigung mit diesem Gegenstande wird aber schon in einer viel früheren Zeit begonnen haben; nach Mitteilungen über eine mündliche Äußerung von \textsc{Gauss}, scheint er im Jahre 1794 die Beziehungen zwischen den arithmetisch-geometrischen Mitteln und den Potenzreihen, in denen die Exponenten mit den Quadratzahlen fortschreiten, gekannt zu haben.

Die Formeln für den in Art. 18. Seite 389. aufgenommenen Algorithmus, der die von \textsc{Gauss} eingeführten neuen Transzendenten mit den Quadratwerten der beiden Argumente zurückführt auf die Transzendenten mit den einfachen Argumentwerten, folgen in einem Handbuch unmittelbar auf eine astronomische Rechnung an deren Schlusse steht `\textit{geendigt d. 2. May} 1809'. Die Aufzeichnungen der anderen Untersuchungen über das Arithmetisch-Geometrische Mittel befinden sich teils auf einzelnen nicht datierten Blättern, teils erscheinen sie in den Handbüchern wegen ihrer Kürze an einigen, früher zu größerer Übersichtlichkeit zwischen verschiedenartigen Arbeiten leer gelassenen, Stellen niedergeschrieben und erlauben keine sichere Zeitangabe.

Wie schon in Art. 12 bemerkt, habe ich geglaubt zur Annehmlichkeit für den Leser diese sehr zerstückelten Untersuchungen durch eine zusammenhängende Darstellung vereinigen zu müssen selbst auf die Gefahr hin, hier einige Entwicklungen hinzustellen, die von \textsc{Gauss} nicht ausgeführt worden sind, wie z. B. die Ableitung der Differentialgleichung für das Arithmetisch-Geometrische Mittel ohne die Reihenentwicklung und die Darstellung durch bestimmte Integrale vorauszusetzen, eine Ableitung, die sich an die in Art. 10. ausgeführte Untersuchung anschließt und sich von der von Herrn \textsc{Borchardt} gegebenen ersten derartigen Ableitung unterscheidet. Die in Art. 17. angeregte Frage über den Zusammenhang zwischen den binären quadratischen Formen mit negativen Determinanten und den von \textsc{Gauss} gefundenen neuen Transzendenten findet ihre vollständigste Erledigung durch die Untersuchungen des Herrn \textsc{Kronecker} über diesen Gegenstand.

\medskip
Für die \textit{lemniskatischen Funktionen} besitzen wir die von \textsc{Gauss} in einem Handbuch verzeichnete Zeitbestimmung `\textit{Wir hatten begonnen, lemniskatische Funktionen zu betrachten} 1797. \textit{Januar.} 8.' Von den im Nachlass vorhandenen Aufzeichnungen scheint nach Papier und Form der Handschrift zu urteilen, die auf einem besonderen Blatte stehenden und hier von mir mit 1, 2, 3, 4 bezeichneten Artikel der \textit{ersten Gruppe} der Untersuchungen über die lemniskatischen Funktionen die früheste zu sein. Die folgenden hier auf Seite 406--412 unter der gemeinsamen von \textsc{Gauss} an mehreren Stellen gebrauchten Überschrift: `\textit{Einige neue Formeln die Lemniscatischen Functionen betreffend}' zusammengestellten Aufzeichnungen derselben ersten Gruppe gehören einer ungleich späteren Zeit an, die darin enthaltenen Resultate sind auch wohl viel früher gefunden und teils zur Gedächtnisprobe, teils in dem Streben recht elegante Formeln zu erhalten von Neuem niedergeschrieben und zwar in zwei Handbüchern, von denen das eine im `\textit{November} 1801' das andere `\textit{im Mai} 1809 \textit{angefangen}' ist. Die Funktionen sin. lemn. sowie cos. lemn. bezeichnet \textsc{Gauss} überall durch die aus der Zusammenziehung von \(s\) und \(l\) sowie von \(c\) und \(l\) gebildeten Schriftzüge. Da diese bis jetzt nicht in Druckzeichen vorhanden sind, so habe ich sie hier durch die Worte selbst ersetzt.

\SourceLeaf{503}

Von den Untersuchungen, welche sich vorzugsweise auf die Darstellung der lemniskatischen Funktionen durch unendliche Produkte und durch trigonometrische Reihen beziehen und welche ich als eine \textit{zweite Gruppe} zusammengestellt habe, bilden die ersten vier Artikel die \textit{Scheda prima} eines nach der Angabe des Titelblattes im Juli 1798 begonnenen Notizbuchs. Die Artikel 1, 2, 3 sind von \textsc{Gauss} selbst nummeriert, dann folgt im selben Heft unmittelbar der Inhalt von Art. [4.] [5.] Seite 415, und hiernach mit vielfachen Unterbrechungen durch astronomische Untersuchungen der Teil der in Art. [17.] Seite 431 aufgenommenen Rechnungen, der mit lateinischem Text erläutert ist, und die einzelnen Teile des Inhalts von Art. [6.] [7.] Seite 417,418. Der Inhalt von Art. [8.] [9.] [10.] Seite 418, 419, Art. [15.] [16.] Seite 423--425 findet sich zerstreut in einem 1799 November angefangenen Notizbuch. Die Aufzeichnungen der übrigen Untersuchungen gehören, vielleicht mit Ausschluss der Fünfteilung des ganzen lemniskatischen Bogens Art. [13.], wohl einer späteren Zeit an und befinden sich, außer den in Art. [14.] wiedergegebenen und in ein Handbuch nach dem 20. Febr. 1817 eingetragenen Summationsformeln für das lemniskatische Integral der zweiten Art, alle auf einzelnen Blättern.

\medskip
Die der Zeit nach erste unter den im handschriftlichen Nachlass erhaltenen die allgemeinen elliptischen Funktionen betreffenden Aufzeichnungen ist wohl die im November 1799 begonnene auf Seite 433--435 abgedruckte Abhandlung \textit{Zur Theorie der neuen Transscendenten} I. die sich in einem Notizbuch dessen Titelblatt die Aufschrift `\textit{Varia, imprimis de Integrali} \(\displaystyle\int\frac{\mathrm{d}u}{\surd(1+\mu\mu\sin u^{2})}\), \textit{Novembr.} 1799' trägt, zwischen Untersuchungen über ganz verschiedenartige Gegenstände zerstreut befindet.

\medskip
Die Abhandlung II. mit der Überschrift `\textit{Zur Theorie der transscendenten Functionen gehörig}' folgt in einem Handbuch unmittelbar nach den in der \textit{Theoria motus corporum coelestium} wiedergegebenen Hilfstafeln. Diese Untersuchungen, insbesondere die auf die Siebenteilung bezüglichen, werden wohl dem Jahre 1808 angehören und die Veranlassung zu einer Mitteilung an \textsc{Schumacher} gewesen sein.

\medskip
Die Abhandlung III. mit den Eingangsworten `Die Theoreme in Bezug auf diejenigen Reihen und unendlichen Produkte, welche zu der Theorie der Arithmetisch-Geometrischen Mittel gehören, ordnen wir so:' Seite 446--460 folgt in einem Handbuch unmittelbar nach einer astronomischen Rechnung, der die Bemerkung beigefügt ist `geendigt den 28. April 1809'.

\medskip
Die Abhandlung IV. `Hundert Theoreme über die neue Transscendente' Seite 461--469 steht auf einzelnen Blättern ohne irgendeine Zeitangabe.

\medskip
Die Abhandlung V. Seite 470--480 mit den Eingangsworten `Allgemeines Theorem' enthält die beiden Zeitangaben 1827 Aug. 6. und Aug. 29. Der letzte Teil dieses Aufsatzes bildet wohl eine der frühesten Vorarbeiten von \textsc{Gauss} für seine `Allgemeinen Lehrsätze in Beziehung auf die im verkehrten Verhältnisse des Quadrats der Entfernung wirkenden Anziehungs- und Abstossungskräfte' und lässt vermuten, dass er den Zusammenhang dieses Gebietes mit einem anderen Gebiete der Analysis, nämlich der \SourceLeaf{504}Theorie der Funktionen mit komplexen Argumenten erkannt habe, ein Zusammenhang welcher \textsc{Riemann} auf ein so fruchtbares Gebiet der Forschung geführt hat.

\medskip
Die Untersuchung des \textit{Pentagramma mirificum}, eines sphärischen Fünfecks, dessen fünf Diagonalen Quadranten sind, befindet sich an zwei getrennten Stellen in einem Handbuch, auf einem besonderen Blatte der Art.\,[4.]. Das, was ich hier als Art.\,[1.] und [2.] bezeichnet habe, ist vor dem 23. Januar 1836 geschrieben, das andere enthält in Art.\,[7.] die Zeitangabe 1843 April 20.

\medskip
Zwischen den unter Nr.\ I. zusammengestellten Untersuchungen über die neuen Transzendenten befindet sich auch der Anfang einer Abhandlung mit der Überschrift ‘\textit{Motus solidi a nullis viribus sollicitati}’. Das Problem ist dort bis zu dem bei der bekannten Auflösung auftretenden elliptischen Integral geführt, sodass zu vermuten steht, \textsc{Gauss} habe erkannt, dass die Gleichungen zur Bestimmung dieser Bewegung mit Hilfe der neuen Transzendenten in endlicher Form erscheinen.

Die Aufzeichnung der hier unter V. zusammengestellten Untersuchungen über die neuen Transzendenten, mit den Zeitangaben 1827 August 6 und August 29 ist wohl durch die \textsc{Jacobi}schen Briefe an \textsc{Schumacher} datiert aus Königsberg von 1827 Juni 13 und Aug.\ 2 deren ersterer die algebraischen Gleichungen für die Dreiteilung und Fünfteilung elliptischer Integrale der andere die Gleichung zwischen den trigonometrischen Tangenten der Argumente für die Transformation beliebigen Grades gibt, veranlasst worden. Diese beiden Briefe sind in der im Monat September 1827 ausgegebenen Nr.\ 123 der ‘Astronomischen Nachrichten’ veröffentlicht, aber im Original zuvor an \textsc{Gauss} mitgeteilt worden, wie aus den \textsc{Gauss}schen Briefen an \textsc{Schumacher} vom 4. und 19. Aug.\ 1827 hervorgeht. Die Beweise jener \textsc{Jacobi}schen Lehrsätze hat dieser selbst durch einen aus Königsberg vom 13. November 1827 datierten in der im Monat Dezember 1827 ausgegebenen Nr.\ 127 der ‘Astronomischen Nachrichten’ abgedruckten Briefe an \textsc{Schumacher} veröffentlicht.

Zu dieser Zeit war noch nicht, aber doch im selben Jahre, durch das zweite Heft des zweiten Bandes des \textsc{Crelle}schen Journals für reine und angewandte Mathematik die Abhandlung von \textsc{Abel} ‘\textit{Recherches sur les fonctions elliptiques} \S.\ I.\,---\,\S.\ VII.’ erschienen, und dieses war wohl diejenige Arbeit \textsc{Abel}s, von der \textsc{Gauss} am 30. Mai 1828 an \textsc{Schumacher} schreibt, ‘die, Ihnen gesagt, mir von meinen eigenen Untersuchungen wohl \(\frac{1}{3}\) vorweggenommen hat, und mit diesen zum Teil selbst bis auf die gewählten bezeichnenden Buchstaben übereinstimmt’.

Auf dieselbe Arbeit bezieht sich wohl die folgende Stelle eines Briefs von \textsc{Crelle} an \textsc{Abel} vom 15. Mai 1828: --- --- Folgendes schreibt mir Herr \textsc{Gauss} aus Göttingen, den ich ebenfalls gebeten hatte, mir etwas über die elliptischen Funktionen zu schicken, mit denen er sich, wie ich erfahren habe, seit mehr als 30 Jahren beschäftigt. ‘Andere Beschäftigungen hindern mich im Augenblick daran, diese Untersuchungen niederzuschreiben. Herr \textsc{Abel} ist mir um mindestens ein Drittel zuvorgekommen. Er hat gerade genau denselben Weg eingeschlagen, von dem ich 1798 ausgegangen bin. Daher wundert es mich keineswegs, dass er zum größten Teil zu denselben Ergebnissen gelangt ist. Da er außerdem in seiner Herleitung so viel Scharfsinn, Durchdringung und Eleganz gezeigt hat, halte ich mich dadurch von der Niederschrift meiner eigenen Untersuchungen für entbunden.’ --- ---

\SourceLeaf{505}

Auch über \textsc{Legendre} besitzen wir einen Ausspruch von \textsc{Gauss}. In einem Briefe ohne Datum schreibt er an \textsc{Olbers} ‘Sie verlangten in Ihrem letzten Briefe [wahrscheinlich derjenige ‘Bremen d.\ 16. Aug.\ 1817, \textit{Empfangen den 25. Aug.}’ bezeichnete; die Briefe aus jenen Monaten sind: O. an G. Juli 17. --- G. an O. Aug.\ 2. --- G. an O. der hier im Auszuge mitgeteilte ohne Datum --- O. an G. Nov.\ 2. --- G. an O. Dec.\ 2.], allerteuerster Freund, mein Urteil über \textsc{Mossotti}s in den Mailänder Ephemeriden gegebene Methode die Bahnen von H. K. zu berechnen. Als ich Ihnen neulich schrieb, war mir der Gegenstand nicht gegenwärtig genug, ob ich gleich jenen Aufsatz früher so weit gelesen hatte, dass ich ein Urteil darüber vorläufig gefasst hatte. In jenem Augenblicke erlaubte mir meine Zeit nicht, mich gleich wieder gehörig in die Sache hineinzustudieren, und ich überging daher Ihre Anfrage. Seitdem habe ich nun wieder Anlass genommen, jenen Aufsatz noch einmal zu lesen, und in den eigentlichen Geist weiter einzudringen, und ich will heute eine Stunde dazu anwenden mich mit Ihnen über diesen Gegenstand zu unterhalten.’

‘Geneigt, wie ich von jeher gewesen bin, jeden neuen originellen oder genialen Gedanken mit Liebe aufzunehmen\HistoricalNote{31}, wurde ich von der wirklich neuen Idee in \textsc{Mossotti}s Aufsatz bei meiner ersten Lektüre frappiert.’ --- ---

\medskip


\medskip
Es verdient noch besonders ausgesprochen zu werden, dass in Bezug auf die Theorie der Teilung des Lemniskatenbogens der handschriftliche Nachlass nichts enthält als die Hilfssätze, die der vorliegende Abdruck dazu bietet, während in dem Werk ‘\textit{Disquisitiones arithmeticae}’, das im Juli 1801 erschienen ist, Art.\ 335 der Sectio septima, de aequationibus circuli sectiones definientibus, gesagt wird --- ‘Im Übrigen reichen die Grundsätze der Theorie, deren Darstellung wir in Angriff nehmen, viel weiter, als sie hier ausgeführt werden. Denn sie lassen sich nicht nur auf die Kreisfunktionen, sondern mit gleichem Erfolg auf viele andere transzendente Funktionen anwenden, zum Beispiel auf diejenigen, die von dem Integral \(\displaystyle\int\frac{\mathrm{d}x}{\surd(1-x^{4})}\) abhängen, und außerdem auch auf verschiedene Arten von Kongruenzen. Da wir aber über jene transzendenten Funktionen ein umfangreiches eigenes Werk vorbereiten und die Kongruenzen in der Fortsetzung der arithmetischen Untersuchungen ausführlich behandelt werden sollen, erschien es angemessen, hier nur die Kreisfunktionen zu betrachten.’ ---

\textsc{Abel} und \textsc{Jacobi} haben \textsc{Gauss}’ Untersuchungen über die elliptischen Funktionen nicht vorgefunden, sie mussten dieses Gebiet der Wissenschaft von Neuem entdecken.

Die speziellen Beziehungen zwischen den Arbeiten von \textsc{Gauss} in diesem Gebiete der reinen Analysis und den Arbeiten von Anderen werde ich in einer besonderen Schrift im Zusammenhange mit einer Geschichte der gesamten wissenschaftlichen Tätigkeit von \textsc{Gauss} darzustellen versuchen, während ich in diesen seinen eigenen Werken angeschlossenen Bemerkungen nur die betreffenden aktenmäßigen Tatsachen aufgenommen habe.

\par Göttingen im Juni 1865.
\begin{flushright}
\textsc{Schering}.
\end{flushright}
\begin{center}
\rule{0.27\linewidth}{0.2pt}
\end{center}

\clearpage
\endgroup
